\section{Conclusion}
\label{sec:conclusion}

This paper addresses the fragility of multimodal visual object tracking under missing-modality conditions by proposing FlexTrack-V2. Extending the adaptive-capacity principle of our conference work, we shift the robustness paradigm from \emph{passive coping} to \emph{active compensation}. The core idea is that an intact modality carries strong priors about an absent one. To exploit this, we introduce the Bilateral Modality-specific feature Reconstruction (BMR) module, which couples with a Heterogeneous Mixture-of-Experts (HMoE) to actively hallucinate missing features while scaling computation dynamically. To train this generative pathway stably without rigid external teachers, we propose Curriculum Missing Augmentation (CMA) with online self-distillation. Extensive experiments across nine benchmarks show that FlexTrack-V2 matches the state-of-the-art on complete modalities and yields substantial, consistent improvements when modalities drop out. 

A current limitation of our framework is that bilateral hallucination assumes the presence of at least one informative modality; if both sensors fail simultaneously (e.g., total occlusion in the dark), the active compensation degrades. Extending the architecture with temporal memory banks to hallucinate features across both time and modalities remains an important direction for future work. Ultimately, we hope FlexTrack-V2 provides a reliable and unified blueprint for deploying multimodal trackers in unpredictable real-world environments.
